\section{行编辑器：\texttt{shell\_readline}}

\codefile{src/kernel/core/shell.c}
\begin{lstlisting}[style=cstyle]
static size_t shell_readline(char* buf, size_t cap) {
    size_t len = 0;
    for (;;) {
        char c = console_getc();  // 阻塞等待一个字符

        if (c == '\n') {
            printk("\n");
            buf[len] = '\0';
            return len;
        }

        // Backspace (0x08) 或 DEL (0x7F)
        if (c == '\b' || (unsigned char)c == 0x7F) {
            if (len > 0) {
                len--;
                printk("\b \b");  // 退格 + 空格覆盖 + 再退格
            }
            continue;
        }

        // 忽略其他控制字符（< 0x20）
        if ((unsigned char)c < 0x20) continue;

        // 可见字符：存入缓冲区 + 回显
        if (len + 1 < cap) {
            buf[len++] = c;
            printk("%c", c);
        }
    }
}
\end{lstlisting}

退格的三步序列 \texttt{\textbackslash b \textbackslash{} \textbackslash b} 详解：

\begin{figure}[H]
  \centering
  % figures/ch06/fig04-textbackslash-b-textbackslash-.tex — auto-extracted from chapter source
\begin{tikzpicture}[node distance=0.5cm]
  \node[font=\ttfamily] (s0) {szy-kernel > helo};
  \node[font=\ttfamily, below=0.5cm of s0] (s1) {szy-kernel > hel\textvisiblespace};
  \node[font=\ttfamily, below=0.5cm of s1] (s2) {szy-kernel > hel};
  \node[right=1cm of s0, font=\scriptsize\itshape] {按 Backspace 前};
  \node[right=1cm of s1, font=\scriptsize\itshape] {\textbackslash b 退格 → 空格覆盖 'o'};
  \node[right=1cm of s2, font=\scriptsize\itshape] {再 \textbackslash b 退格到正确位置};
\end{tikzpicture}

  \caption{退格三步序列：\texttt{\textbackslash b} (退格) → \texttt{' '} (覆盖) → \texttt{\textbackslash b} (归位)}
\end{figure}

\subsection{命令分词：\texttt{shell\_tokenize}}

\codefile{src/kernel/core/shell.c}
\begin{lstlisting}[style=cstyle]
static int shell_tokenize(char* line, char** argv, int argv_cap) {
    int argc = 0;
    char* p = line;
    while (*p) {
        while (*p && is_space(*p)) p++;  // 跳过空白
        if (!*p) break;
        if (argc >= argv_cap) break;
        argv[argc++] = p;               // 记录 token 起始
        while (*p && !is_space(*p)) p++; // 跳过非空白
        if (*p) { *p = '\0'; p++; }     // 用 NUL 截断 token
    }
    return argc;
}
\end{lstlisting}

\begin{figure}[H]
  \centering
  % figures/ch06/fig05-tokenize-textbackslash-0-argv-.tex — auto-extracted from chapter source
\begin{tikzpicture}[node distance=0cm]
  % 输入行
  \foreach \i/\ch in {0/p, 1/m, 2/m, 3/\textvisiblespace, 4/a, 5/l, 6/l, 7/o, 8/c, 9/\textvisiblespace, 10/3, 11/\textbackslash 0} {
    \node[memcell, minimum width=0.9cm, minimum height=0.6cm, font=\ttfamily\small] (c\i) at (\i*0.95, 0) {\ch};
  }
  \node[font=\small, above=0.4cm of c5] {\texttt{line[]} 内容};

  % Tokenize 后
  \foreach \i/\ch in {0/p, 1/m, 2/m, 3/\textbackslash 0, 4/a, 5/l, 6/l, 7/o, 8/c, 9/\textbackslash 0, 10/3, 11/\textbackslash 0} {
    \node[memcell, minimum width=0.9cm, minimum height=0.6cm, font=\ttfamily\small] (d\i) at (\i*0.95, -1.5) {\ch};
  }
  \node[font=\small, above=0.4cm of d5] {tokenize 后（空格变 \textbackslash 0）};

  % argv pointers
  \draw[pointer, red] ([yshift=-0.4cm]d0.south) -- (d0.south) node[below=0.5cm, font=\scriptsize\ttfamily] {argv[0]="pmm"};
  \draw[pointer, blue] ([yshift=-0.4cm]d4.south) -- (d4.south) node[below=0.5cm, font=\scriptsize\ttfamily] {argv[1]="alloc"};
  \draw[pointer, green!50!black] ([yshift=-0.4cm]d10.south) -- (d10.south) node[below=0.5cm, font=\scriptsize\ttfamily] {argv[2]="3"};
\end{tikzpicture}

  \caption{tokenize 将空格替换为 \texttt{\textbackslash 0}，argv[] 指向各 token 起始}
\end{figure}

\begin{warnbox}
\texttt{shell\_tokenize} 会\textbf{修改}输入字符串（用 \texttt{\textbackslash 0} 替换空格）。
这是 C 中常见的"零拷贝分词"技巧——不分配新内存，但调用者不能再使用原始字符串。
\end{warnbox}

% ============================================================================

